\chapter{Independence and conditional expectation}
\label{ch:cond}

\begin{question}
How does partial information change our best guess of a random quantity, and when does information about one quantity tell us nothing about another?
\end{question}

\begin{whatfor}
\begin{itemize}
  \item Conditional probability, total probability and Bayes.
  \item Independence of events, of random variables and of $\sigma$-fields, and what it does to expectations.
  \item The conditional expectation $\E_P(X\given\G)$: a random variable, the best guess of $X$ given the information $\G$.
  \item Its four working rules CE1--CE4, used on every page of Parts~1--4.
\end{itemize}
\skipline{fs:cond}
\end{whatfor}

\begin{prereq}
From Chapters~1--3: $\sigma$-field, atoms, filtration, random variable, $\sigma(\xi)$, density, joint density, expectation $\E_P$, variance, covariance, the normal exponential moment.
\end{prereq}

%=====================================================================
\section{Conditional probability}

\begin{definition}[Conditional probability]
For events $A,B$ with $P(B)\ne0$, the \dfn{conditional probability} of $A$ given $B$ is
\[
  P(A\given B)=\frac{P(A\cap B)}{P(B)} .
\]
\end{definition}

Conditioning on $B$ throws away the outcomes outside $B$ and rescales what is left so that $B$ has probability $1$. Die, told ``even'': the weights become $\frac13$ on each of $2,4,6$, so $P(\{6\}\given\text{even})=\frac{1/6}{1/2}=\frac13$.

\begin{result}[Total probability and Bayes]
\label{res:total}
Let $B_1,B_2,\ldots$ be a \dfn{partition} of $\Omega$: pairwise disjoint events with union $\Omega$, each with $P(B_i)>0$. Then for every event $A$,
\[
  P(A)=\sum_i P(A\given B_i)\,P(B_i),\qquad
  P(B_k\given A)=\frac{P(A\given B_k)\,P(B_k)}{\sum_iP(A\given B_i)\,P(B_i)} .
\]
\end{result}
\begin{proof}
$A$ is the disjoint union of the pieces $A\cap B_i$, so $P(A)=\sum_iP(A\cap B_i)=\sum_iP(A\given B_i)P(B_i)$. Bayes is $P(B_k\given A)=P(A\cap B_k)/P(A)$ with this denominator.
\end{proof}

\begin{example}[Which machine made the defect?]
Machine $M_1$ makes $60\%$ of items, $2\%$ of them defective; $M_2$ makes $40\%$, $5\%$ defective.
\begin{steps}
  \item Total probability: $P(D)=0.02(0.6)+0.05(0.4)=0.012+0.020=0.032$.
  \item Bayes: $P(M_2\given D)=0.020/0.032=0.625$. $M_2$ makes $40\%$ of the items but $62.5\%$ of the defects.
\end{steps}
\end{example}

%=====================================================================
\section{Independence}

\begin{definition}[Independence]
\label{def:indep}
Independence is defined in three layers.
\begin{itemize}
  \item Events $A,B$ are \dfn{independent} if $P(A\cap B)=P(A)P(B)$, i.e.\ $P(A\given B)=P(A)$ when $P(B)>0$. Events $A_1,\dots,A_n$ are independent if $P(A_{i_1}\cap\dots\cap A_{i_k})=P(A_{i_1})\cdots P(A_{i_k})$ for every choice $i_1<\dots<i_k$.
  \item $\sigma$-fields $\G$ and $\mathcal H$ are independent if every $A\in\G$ and $B\in\mathcal H$ are independent events.
  \item Random variables $\xi,\eta$ are independent if $\sigma(\xi)$ and $\sigma(\eta)$ are independent: $P(\xi\in A,\eta\in B)=P(\xi\in A)P(\eta\in B)$ for all Borel $A,B$. A family of random variables is independent if any finitely many of them are. A random variable $\xi$ is independent of $\G$ if $\sigma(\xi)$ and $\G$ are independent.
\end{itemize}
\end{definition}

\begin{example}[Pairwise is not enough]
Two dice. $A$ = first even, $B$ = second even, $C$ = sum even; each has probability $\frac12$.
\begin{itemize}
  \item Every pair: $P(A\cap B)=P(A\cap C)=P(B\cap C)=\frac14=\frac12\cdot\frac12$.
  \item All three: $P(A\cap B\cap C)=P(\text{both even})=\frac14\ne\frac18$.
\end{itemize}
Any two of them decide the third, so the three are not independent.
\end{example}

Disjoint events with positive probabilities are never independent: $P(A\cap B)=0<P(A)P(B)$. Learning $A$ rules $B$ out completely.

\begin{result}[Independence and expectations]
\label{res:indep-E}
Let $\xi,\eta$ be independent and $f,g$ Borel functions with $f(\xi),g(\eta)$ integrable. Then
\[
  \E_P\big[f(\xi)\,g(\eta)\big]=\E_P f(\xi)\;\E_P g(\eta).
\]
In particular $\Cov(\xi,\eta)=0$ and $\Var(\xi+\eta)=\Var\xi+\Var\eta$. If $\xi,\eta$ have densities, they are independent exactly when the joint density factorises: $f_{\xi,\eta}(x,y)=f_\xi(x)f_\eta(y)$.
\end{result}
\begin{proof}
For indicators, $f=\ind{A}$, $g=\ind{B}$, the claim is $P(\xi\in A,\eta\in B)=P(\xi\in A)P(\eta\in B)$, the definition. It extends to step functions by linearity, and to the rest by the limits of Chapter~3.
\end{proof}

\begin{result}[Sums of independent normals]
\label{res:normal-sum}
If $\xi\sim N(\mu_1,\sigma_1^2)$ and $\eta\sim N(\mu_2,\sigma_2^2)$ are independent, then $\xi+\eta\sim N(\mu_1+\mu_2,\ \sigma_1^2+\sigma_2^2)$.
\end{result}
\begin{proof}
Compare exponential moments.
\begin{steps}
  \item By independence and \cref{res:normal-moments}, $\E_Pe^{u(\xi+\eta)}=\E_Pe^{u\xi}\,\E_Pe^{u\eta}=e^{u\mu_1+u^2\sigma_1^2/2}\,e^{u\mu_2+u^2\sigma_2^2/2}$.
  \item This is $e^{u(\mu_1+\mu_2)+u^2(\sigma_1^2+\sigma_2^2)/2}$, the exponential moment of $N(\mu_1+\mu_2,\sigma_1^2+\sigma_2^2)$.
  \item A distribution whose exponential moments $\E_Pe^{u\xi}$ are finite for all $u$ is determined by them (a fact we use without proof). \qedhere
\end{steps}
\end{proof}

%=====================================================================
\section{Conditional expectation}

Start with the simplest information: whether an event $B$ happened.

\begin{definition}[Conditional expectation given an event]
For an event $B$ with $P(B)>0$, $\E_P(X\given B)=\dfrac{1}{P(B)}\displaystyle\int_BX\,dP=\dfrac{\E_P(X\ind B)}{P(B)}$: the average of $X$ over the outcomes in $B$.
\end{definition}

Now suppose the information is a $\sigma$-field $\G\subseteq\F$ with atoms $G_1,\dots,G_k$ (a partition of $\Omega$). Once we learn which atom happened, our best guess of $X$ is its average over that atom. This guess depends on the outcome, so it is a random variable.

\begin{definition}[Conditional expectation given a $\sigma$-field]
\label{def:cond-exp}
Let $X$ be integrable and $\G\subseteq\F$ a $\sigma$-field. The \dfn{conditional expectation} $\E_P(X\given\G)$ is the random variable $Y$ with
\begin{enumerate}[label=(\roman*)]
  \item $Y$ is $\G$-measurable (it uses only the information in $\G$);
  \item $\int_GY\,dP=\int_GX\,dP$ for every $G\in\G$ (it has the right average on every event of $\G$).
\end{enumerate}
Such a $Y$ exists and is unique up to a set of probability $0$. We write $\E_P(X\given\xi)$ for $\E_P(X\given\sigma(\xi))$, and $\E_P(X\given\xi=y)$ for its value where $\xi=y$.
\end{definition}

\begin{result}[The partition formula]
\label{res:partition}
If $\G$ has atoms $G_1,\dots,G_k$ with $P(G_i)>0$, then
\[
  \E_P(X\given\G)=\sum_{i=1}^k\E_P(X\given G_i)\,\ind{G_i}:
\]
on each atom, it equals the average of $X$ over that atom.
\end{result}
\begin{proof}
(i) forces $Y$ to be constant on each atom, say $c_i$ on $G_i$. Then (ii) with $G=G_i$ reads $c_iP(G_i)=\int_{G_i}X\,dP$, so $c_i=\E_P(X\given G_i)$.
\end{proof}

\begin{example}[A die told odd or even]
\label{ex:die-ce}
$X(\omega)=\omega$, $\G=\{\varnothing,\{1,3,5\},\{2,4,6\},\Omega\}$.
\begin{steps}
  \item Odd atom: $\E_P(X\given\text{odd})=\frac{1+3+5}{3}=3$. Even atom: $\frac{2+4+6}{3}=4$.
  \item So $\E_P(X\given\G)=3$ on odd outcomes and $4$ on even ones: a random variable with two values.
  \item Check the average: $\frac12\cdot3+\frac12\cdot4=3.5=\E_P(X)$.
\end{steps}
\end{example}

\begin{pic}
\begin{tikzpicture}
\begin{axis}[width=8.5cm,height=4.4cm,ybar,bar width=12pt,ymin=0,ymax=6.8,xtick={1,...,6},xlabel={outcome $\omega$},ylabel={value},xmin=0.4,xmax=6.6]
  \addplot[fill=prec!25,draw=prec] coordinates {(1,1)(2,2)(3,3)(4,4)(5,5)(6,6)};
  \addplot[ybar=0pt,draw=none,forget plot] coordinates {(1,0)};
\end{axis}
\begin{axis}[width=8.5cm,height=4.4cm,ymin=0,ymax=6.8,xmin=0.4,xmax=6.6,axis lines=none]
  \addplot[thmc,very thick,mark=*] coordinates {(1,3)(3,3)(5,3)};
  \addplot[defc,very thick,mark=*] coordinates {(2,4)(4,4)(6,4)};
  \node[thmc,font=\scriptsize,anchor=east] at (axis cs:0.95,3) {odd: $3$};
  \node[defc,font=\scriptsize,anchor=east] at (axis cs:1.9,4.4) {even: $4$};
\end{axis}
\end{tikzpicture}
\picnote{Grey bars: $X(\omega)=\omega$. Lines: $\E_P(X\given\G)$ replaces $X$ by its average on each atom of $\G$: all odd outcomes get $3$, all even outcomes get $4$. Finer information means finer averaging.}
\end{pic}

\begin{result}[With a joint density]
\label{res:cond-density}
If $(X,Y)$ has joint density $f(x,y)$ and $f_Y(y)>0$, the \dfn{conditional density} of $X$ given $Y=y$ is $f_{X\given Y=y}(x)=\dfrac{f(x,y)}{f_Y(y)}$, and
\[
  \begin{aligned}
  \E_P(X\given Y=y)&=\int x\,f_{X\given Y=y}(x)\,dx,\\
  P(X>t)&=\int P(X>t\given Y=y)\,f_Y(y)\,dy=\E_P\big[P(X>t\given Y)\big].
  \end{aligned}
\]
\end{result}
\begin{proof}
For the second formula, swap the order of integration:
\[
  \begin{aligned}
  \int P(X>t\given Y=y)f_Y(y)\,dy&=\int\!\!\int_t^\infty\frac{f(x,y)}{f_Y(y)}\,dx\,f_Y(y)\,dy\\
  &=\int_t^\infty\!\!\int f(x,y)\,dy\,dx=\int_t^\infty f_X(x)\,dx=P(X>t). \qedhere
  \end{aligned}
\]
\end{proof}

%=====================================================================
\section{The rules CE1--CE4}

These four rules, plus two extreme cases, are all we use to compute conditional expectations.

\begin{theorem}[Rules of conditional expectation]
\label{thm:ce-rules}
Let $X,Y$ be integrable, $a,b$ numbers, and $\G_0\subseteq\G_1\subseteq\F$ $\sigma$-fields.
\begin{itemize}
  \item[\textbf{CE1}] Linearity: $\E_P(aX+bY\given\G)=a\,\E_P(X\given\G)+b\,\E_P(Y\given\G)$.
  \item[\textbf{CE2}] Taking out what is known: if $Z$ is $\G$-measurable, $\E_P(ZX\given\G)=Z\,\E_P(X\given\G)$.
  \item[\textbf{CE3}] Independence: if $X$ is independent of $\G$, then $\E_P(X\given\G)=\E_P(X)$.
  \item[\textbf{CE4}] Tower: $\E_P\big[\E_P(X\given\G_1)\given\G_0\big]=\E_P(X\given\G_0)$. With $\G_0=\{\varnothing,\Omega\}$: $\E_P\big[\E_P(X\given\G_1)\big]=\E_P(X)$.
\end{itemize}
Two extremes: if $X$ is $\G$-measurable, $\E_P(X\given\G)=X$ (full information); if $\G=\{\varnothing,\Omega\}$, $\E_P(X\given\G)=\E_P(X)$ (no information).
\end{theorem}
\begin{proof}
We prove the rules on a finite $\Omega$, where the partition formula (\cref{res:partition}) gives everything; the general case uses the same two defining properties.
\begin{steps}
  \item CE1: on each atom, the average of $aX+bY$ is $a$ times the average of $X$ plus $b$ times that of $Y$.
  \item CE2: $Z$ is constant on each atom $G$ of $\G$, say $z$, so the average of $ZX$ over $G$ is $z$ times the average of $X$.
  \item CE3: for an atom $G$ of $\G$, independence gives $\E_P(X\ind G)=\E_P(X)P(G)$, so the average over $G$ is $\E_P(X)$.
  \item CE4: let $B$ be an atom of $\G_0$. It is a union of atoms of $\G_1$, so
  \[
    \E_P\big[\E_P(X\given\G_1)\given B\big]=\frac{1}{P(B)}\int_B\E_P(X\given\G_1)\,dP=\frac{1}{P(B)}\int_BX\,dP=\E_P(X\given B),
  \]
  where the middle step is property (ii) of $\E_P(X\given\G_1)$, applied to $B\in\G_1$. \qedhere
\end{steps}
\end{proof}

\begin{intuition}
CE4 says: averaging the finer guess with coarser information gives the coarser guess. The coarse information always wins.
\end{intuition}

\begin{example}[Two dice]
\label{ex:two-dice-ce}
$X_1,X_2$ two independent dice, $\G=\sigma(X_1)$; $\E_P(X_2)=3.5$, $\E_P(X_2^2)=\frac{91}{6}$.
\begin{itemize}
  \item $\E_P(X_1+X_2\given X_1)=X_1+3.5$ (CE1, then known, then CE3).
  \item $\E_P(X_1X_2\given X_1)=X_1\,\E_P(X_2\given X_1)=3.5X_1$ (CE2, CE3). If $X_1=2$ it is $7$.
  \item $\E_P\big((X_1+X_2)^2\given X_1\big)=X_1^2+2X_1(3.5)+\frac{91}{6}$ (expand, then CE1--CE3).
  \item Tower: $\E_P\big[\E_P(X_1+X_2\given X_1)\big]=\E_P(X_1)+3.5=7=\E_P(X_1+X_2)$.
\end{itemize}
\end{example}

%=====================================================================
\formulasheet{fs:cond}

\begin{mastercard}[label=card:bayes]{Conditioning on a partition (total probability and Bayes)}
For a partition $B_1,B_2,\dots$ of $\Omega$ (disjoint, covering, each $P(B_i)>0$):
\[
  \boxed{\begin{gathered}P(B_k\given A)=\frac{P(A\given B_k)\,P(B_k)}{\sum_iP(A\given B_i)\,P(B_i)},\\ P(A\given B)=\frac{P(A\cap B)}{P(B)}.\end{gathered}}
\]
\begin{plugrows}
  \plugrow{Probability of $A$ by cases}{the denominator alone}{$P(A)=\sum_iP(A\given B_i)P(B_i)$}{partition}{machines: $P(D)=0.02(0.6)+0.05(0.4)=0.032$}
  \plugrow{Reverse the conditioning}{the full formula}{$P(B_k\given A)$}{partition}{$P(M_2\given D)=0.020/0.032=0.625$}
  \plugrow{Random parameter, continuous}{partition by the value $y$ of $Y$; sum becomes integral}{$P(A)=\int P(A\given Y=y)f_Y(y)\,dy$}{$Y$ has a density}{rate $1$ or $2$ w.p.\ $\frac12$: $P(X>1)=\frac12e^{-1}+\frac12e^{-2}\approx0.252$}
\end{plugrows}
\end{mastercard}

\begin{mastercard}[label=card:indep]{Independence}
$X_1,\dots,X_n$ are independent exactly when expectations of products split, for all functions $f_i$:
\[
  \boxed{\begin{gathered}\E_P\Big[\prod_{i=1}^nf_i(X_i)\Big]=\prod_{i=1}^n\E_P\,f_i(X_i).\end{gathered}}
\]
\begin{plugrows}
  \plugrow{Events}{$f_i=$ indicators}{$P(A_{i_1}\cap\dots\cap A_{i_k})=P(A_{i_1})\cdots P(A_{i_k})$ for every sub-collection}{all sub-collections, not just pairs}{two dice, first even / second even / sum even: pairwise yes, all three no}
  \plugrow{Densities}{compare the joint density with the product of the marginals}{independent $\iff f(x,y)=f_X(x)f_Y(y)$}{joint density exists}{uniform on $[0,1]^2$: $1=1\cdot1$, independent}
  \plugrow{Covariance and variance}{$f_1=f_2=$ identity}{$\Cov=0$, $\Var(\sum X_i)=\sum\Var X_i$}{$\Cov=0$ does not give independence}{two independent dice: $\Var(X_1+X_2)=2\cdot\frac{35}{12}$}
  \plugrow{Sums of independent normals}{$f_i=e^{ua_ix}$ and Card~\ref{card:normal-mgf}}{$\sum a_iX_i\sim N\big(\sum a_i\mu_i,\ \sum a_i^2\sigma_i^2\big)$}{independent normal $X_i$}{$N(1,2)$, $N(-3,5)$: $2X-Y\sim N(5,13)$}
  \plugrow{Disjoint events}{$A\cap B=\varnothing$}{never independent if both have positive probability}{$P(A),P(B)>0$}{$P(A\cap B)=0<P(A)P(B)$}
\end{plugrows}
\end{mastercard}

\begin{mastercard}[label=card:freeze]{Conditional expectation: freeze what is known, average the rest}
Let $K$ be known from $\G$ ($\G$-measurable) and $Z$ independent of $\G$. Then the guess given $\G$ is found by treating $K$ as a constant $k$ and averaging over $Z$:
\[
  \boxed{\begin{gathered}\E_P\big(f(K,Z)\given\G\big)=g(K),\quad g(k)=\E_P\,f(k,Z);\\ \E_P\big[\E_P(X\given\G)\given\G_0\big]=\E_P(X\given\G_0)\ \ (\G_0\subseteq\G).\end{gathered}}
\]
\begin{plugrows}
  \plugrow{Linear combinations (CE1)}{$f=ak+bz$}{$\E_P(aK+bZ\given\G)=aK+b\E_PZ$}{always}{two dice: $\E_P(X_1+X_2\given X_1)=X_1+3.5$}
  \plugrow{Take out what is known (CE2)}{$f=k\,z$}{$\E_P(KZ\given\G)=K\E_PZ$}{$K$ known}{$\E_P(X_1X_2\given X_1)=3.5X_1$}
  \plugrow{Independent of the information (CE3)}{$f$ depends on $z$ only}{$\E_P(Z\given\G)=\E_PZ$}{independent of all of $\G$}{$\E_P(X_2\given X_1)=3.5$}
  \plugrow{Future of a sum}{$K=S_n$, $Z=S_{n+m}-S_n$}{$\E_P(S_{n+m}\given\F_n)=S_n+m\mu$; $\E_P(S_{n+m}^2\given\F_n)=(S_n+m\mu)^2+m\sigma^2$}{i.i.d.\ steps, mean $\mu$, variance $\sigma^2$}{fair $\pm1$: $\E_P(S_5^2\given\F_3)=S_3^2+2$}
  \plugrow{Finite information}{$\G$ with atoms $G_i$}{on $G_i$: $\E_P(X\given\G)=\E_P(X\ind{G_i})/P(G_i)$, the average over the atom}{$P(G_i)>0$}{die, odd/even: $3$ on odd, $4$ on even}
  \plugrow{Unconditional mean (tower)}{$\G_0=\{\varnothing,\Omega\}$}{$\E_P[\E_P(X\given\G)]=\E_PX$}{always}{$\frac12\cdot3+\frac12\cdot4=3.5$}
  \plugrow{Given a value, with a density}{$f_{X\given Y=y}=f(x,y)/f_Y(y)$}{$\E_P(X\given Y=y)=\int x\,f_{X\given Y=y}(x)dx$}{$f_Y(y)>0$}{used for Brownian motion in Chapter~9}
\end{plugrows}
\end{mastercard}
